\documentclass[a4paper, 11pt]{article}

\usepackage{amsmath, amssymb, mathtools, mathrsfs}
\usepackage{cite}
\usepackage{hyperref}
\usepackage{cancel}
\usepackage{tikz}
\usepackage{tcolorbox}
\usepackage[normalem]{ulem}
\usepackage[a4paper,margin=1in]{geometry}

\hfuzz=90pt
\numberwithin{equation}{section}

\newcommand{\ket}[1]{| #1 \rangle}
\newcommand{\bra}[1]{\langle #1 |}
\newcommand{\braket}[2]{| #1 \rangle\langle #2 |}

\begin{document}
{\LARGE \textbf{Notes} \par}\vspace{1em}

We setup a system to observe the evolution of $N$ three-level systems,
each arranged in a $\Lambda$ configuration. We denote the ground states as $\ket{g}$
and $\ket{f}$, and the excited state as $\ket{e}$.

We drive the transition $\ket{g}\leftrightarrow\ket{e}$ with a classical laser, and
allow a waveguide to mediate the transitions $\ket{f}\leftrightarrow\ket{e}$ and $\ket{g}\leftrightarrow\ket{e}$.

\vspace{1em}{\bf\color{red} Note:} There are certain equations, or certain terms in
equations that are colored in red. These are terms that have been thrown out due to
some assumption that has been made. They have not been removed entirely in case said
assumption is not valid.

\section{Hamiltonian}

We build up the Hamiltonian from four distinct contributions:
\begin{equation}
  \hat{H} = \hat{H}_F + \hat{H}_A + \hat{H}_{AF} + \hat{H}_{drive}\\
\end{equation}

Here, $\hat{H}_F$ is the free-field Hamiltonian, $\hat{H}_A$ is the atomic
Hamiltonian, $\hat{H}_{AF}$ is the atom-field interaction, and $\hat{H}_{drive}$ is an
external drive.

We can now define the individual components:
\begin{align*}
  \hat{H}_F &= \int d^3r \int_0^\infty d\omega \hbar\omega \hat{\mathbf{f}}^\dagger
  (\mathbf{r}, \omega) \hat{\mathbf{f}}(\mathbf{r}, \omega)\\
  \hat{H}_A &= \sum_{i=1}^N \hbar\omega_e \hat{\sigma}_{ee}^i + \hbar\omega_f
  \hat{\sigma}_{ff}^i\\
  \hat{H}_{AF} &= -\sum_{i=1}^N \hat{\mathbf{E}}(\mathbf{r}_i, t) \cdot \left(
  \mathbf{d}_i \hat{\sigma}_i + \mathbf{d}_i^* \hat{\sigma}_i^\dagger \right)\\
  \hat{H}_{drive} &= -\sum_{i=1}^N \hbar \left( \frac{\Omega}{2} e^{i\omega_Lt}
  \hat{\sigma}_i + \frac{\Omega^*}{2} e^{-i\omega_Lt} \hat{\sigma}_i^\dagger \right)
\end{align*}

The ground state $\ket{g}$ is taken as the reference with zero energy. Throughout, we
use the notation $$\hat{\sigma}_{ab}^i \equiv \braket{a_i}{b_i}$$ to denote the atomic
operator on the $i$-th atom.
\vspace{1em}

In the case of our system, we can expand out the Hamiltonians for the atom-field
interaction as well as the external drive, giving us this total Hamiltonian
($\hbar=1$):

\begin{equation}
  \begin{aligned}
    \hat{H} = &\int d^3r \int_0^\infty \omega \hat{\mathbf{f}}^\dagger (\mathbf{r},
    \omega) \hat{\mathbf{f}}(\mathbf{r}, \omega) d\omega + \sum_{i=1}^N \omega_e
    \hat{\sigma}_{ee}^i + \omega_f \hat{\sigma}_{ff}^i\\ &-\sum_{i=1}^N
    \hat{\mathbf{E}}(\mathbf{r}_i, t) \cdot \left( \mathbf{d}_{ge} \hat{\sigma}^i_{ge}
    + \mathbf{d}_{ge}^* \hat{\sigma}^i_{eg} + \mathbf{d}_{fe} \hat{\sigma}^i_{fe} +
    \mathbf{d}_{fe}^* \hat{\sigma}^i_{ef} \right)\\ &-\sum_{i=1}^N \frac{\Omega}{2}
    e^{i\omega_Lt} \hat{\sigma}_{ge}^i + \frac{\Omega^*}{2} e^{-i\omega_Lt}
    \hat{\sigma}_{eg}^i
  \end{aligned}
\end{equation}

\subsection{Rotating Frame}
We rotate out the laser frequency using the unitary transformation 
\begin{equation}
  \hat{U} = \sum_{i=1}^N \exp\left(it\left[\omega_L\,\hat{\sigma}_{ee}^i +
      \omega_f\hat{\sigma}_{ff}^i\right]\right)
      % + \int d^3r \int_0^\infty \omega
      % \hat{\mathbf{f}}^\dagger (\mathbf{r}, \omega) \hat{\mathbf{f}}(\mathbf{r},
      % \omega) d\omega\right]\right)
\end{equation}


The Hamiltonian evolves by \cite{SteckQAO}
\begin{equation}
  \hat{\tilde{H}} = \hat{U}\hat{H}\hat{U}^\dagger + i(\partial_t \hat{U})\hat{U}^\dagger
\end{equation}

After defining the detuning $\Delta = \omega_e - \omega_L$, we get the transformed
Hamiltonian:
\begin{equation}
  \begin{aligned}
    \hat{\tilde{H}} = & \int d^3r \int_0^\infty \omega \hat{\mathbf{f}}^\dagger
    (\mathbf{r}, \omega) \hat{\mathbf{f}}(\mathbf{r}, \omega) d\omega + \sum_{i=1}^N
    \Delta \hat{\sigma}_{ee}^i -\sum_{i=1}^N \hat{\mathbf{E}}(\mathbf{r}_i, t) \cdot
    \\&\left( e^{-i\omega_Lt}\mathbf{d}_{ge} \hat{\sigma}^i_{ge} +
    e^{i\omega_Lt} \mathbf{d}_{ge}^* \hat{\sigma}^i_{eg} + e^{-i(\omega_L - \omega_f)t}
    \mathbf{d}_{fe} \hat{\sigma}^i_{fe} + e^{i(\omega_L-\omega_f)t} \mathbf{d}_{fe}^*
    \hat{\sigma}^i_{ef} \right)\\ &-\sum_{i=1}^N \left(\frac{\Omega}{2}
    \hat{\sigma}_{ge}^i + \frac{\Omega^*}{2} \hat{\sigma}_{eg}^i\right)
  \end{aligned}
\end{equation}

We also define the new transformed operators
\begin{align}
  \hat{E}_{g} &= \hat{\mathbf{E}}(\mathbf{r}_i, t) \cdot \mathbf{d}_{ge} \cdot
  e^{-i\omega_Lt}\\
  \hat{E}^*_{g} &= \hat{\mathbf{E}}(\mathbf{r}_i, t) \cdot \mathbf{d}^*_{ge} \cdot
  e^{i\omega_Lt}\\
  \hat{E}_{f} &= \hat{\mathbf{E}}(\mathbf{r}_i, t) \cdot \mathbf{d}_{fe} \cdot
  e^{-i(\omega_L-\omega_f)t}\\
  \hat{E}^*_{f} &= \hat{\mathbf{E}}(\mathbf{r}_i, t) \cdot \mathbf{d}^*_{fe} \cdot
  e^{i(\omega_L-\omega_f)t}
\end{align}

Using this, we rewrite out transformed Hamiltonian:
\begin{equation}
  \begin{aligned}
    \hat{H} = & \int d^3r \int_0^\infty \omega \hat{\mathbf{f}}^\dagger (\mathbf{r},
    \omega) \hat{\mathbf{f}}(\mathbf{r}, \omega) d\omega + \sum_{i=1}^N \left(\Delta
      \hat{\sigma}_{ee}^i + \omega_f
    \hat{\sigma}_{ff}^i\right)\\&-\sum_{i=1}^N \left(\hat{E}_g\hat{\sigma}_{ge}^i +
    \hat{E}^*_g\hat{\sigma}_{eg}^i + \hat{E}_f\hat{\sigma}_{fe}^i + \hat{E}^*_f
    \hat{\sigma}_{ef}^i\right)- \left( \frac{\Omega}{2} \hat{\sigma}_{ge}^i +
    \frac{\Omega^*}{2}\hat{\sigma}_{eg}^i\right)
  \end{aligned}
\end{equation}

\subsection{Adiabatic Elimination}
Following the operator formalism by Reiter and S\o{}rensen \cite{PhysRevA.85.032111},
we split up the Hamiltonian into ground and excited state subspaces as well as the
excitations and deexcitations:
\begin{equation}
  \hat{H} = \hat{H}_g + \hat{H}_e + \hat{V}_+ + \hat{V}_-
\end{equation}

Below, we have the separated components of the Hamiltonian
\begin{align*}
  \hat{H}_g &= \int d^3r \int_0^\infty \omega \hat{\mathbf{f}}^\dagger (\mathbf{r},
  \omega) \hat{\mathbf{f}}(\mathbf{r}, \omega) d\omega + \sum_{i=1}^N\omega_f
  \hat{\sigma}_{ff}^i\\
  \hat{H}_e &= \sum_{i=1}^N \Delta \hat{\sigma}_{ee}^i\\
  \hat{V}_+ &= -\sum_{i=1}^N \left[\left( \hat{E}_g^* + \frac{\Omega^*}{2}\right)
  \hat{\sigma}_{eg}^i + \hat{E}_f^* \hat{\sigma}_{ef}^i\right]\\
  \hat{V}_- &= -\sum_{i=1}^N \left[\left( \hat{E}_g + \frac{\Omega}{2}\right)
  \hat{\sigma}_{ge}^i + \hat{E}_f \hat{\sigma}_{fe}^i\right]
\end{align*}

The effective eliminated Hamiltonian is then defined
\begin{equation}
  \hat{H}_{\text{eff}} = -V_-  \hat{H}_e^{-1}  V_+ + \hat{H}_g
\end{equation}

Due to the selection rules, we end up with an effective Hamiltonian
\begin{align*}
  \hat{H} &=-\frac{1}{\Delta} \sum_{i=1}^N \left[\left( \hat{E}_g +
  \frac{\Omega}{2}\right) \hat{\sigma}_{ge}^i + \hat{E}_f \hat{\sigma}_{fe}^i\right]
  \hat{\sigma}_{ee}^i \left[\left( \hat{E}_g^* + \frac{\Omega^*}{2}\right)
  \hat{\sigma}_{eg}^i + \hat{E}_f^* \hat{\sigma}_{ef}^i\right] + \hat{H}_g\\
  &= -\frac{1}{\Delta} \sum_{i=1}^N \left[ \left( \hat{E}_g + \frac{\Omega}{2}\right)
  \left( \hat{E}_g^* + \frac{\Omega^*}{2}\right)\hat{\sigma}_{gg}^i +
  \left( \hat{E}_g + \frac{\Omega}{2}\right) \hat{E}_f^*\hat{\sigma}_{gf}^i\right.\\
  &\hspace{5em}+ \left.\hat{E}_f \left( \hat{E}_g^* +
  \frac{\Omega^*}{2}\right)\hat{\sigma}_{fg}^i + \hat{E}_f
  \hat{E}_f^*\hat{\sigma}_{ff}^i\right] + \hat{H}_g\\
\end{align*}

Here, we make the assumption that the drive is entirely real, i.e. $\Omega =
\Omega^*$. Now, the adiabatically eliminated Hamiltonian is

\begin{equation}
  \begin{aligned}
    \hat{H} = & \int d^3r \int_0^\infty \omega \hat{\mathbf{f}}^\dagger (\mathbf{r},
    \omega) \hat{\mathbf{f}}(\mathbf{r}, \omega) d\omega + \sum_{i=1}^N \left(
    \frac{\hat{E}_f \hat{E}_f^*}{\Delta} + \omega_f\right)\hat{\sigma}_{ff}^i\\
    &+\frac{1}{\Delta} \sum_{i=1}^N \left(\hat{E}_g \hat{E}_g^* + \frac{\hat{E}_g
    \cdot \Omega}{2} + \frac{\Omega \cdot \hat{E}^*_g}{2} + \frac{\Omega^2}{4}\right)
    \hat{\sigma}_{gg}^i \\ &+ \frac{1}{\Delta}\sum_{i=1}^N \left[ \left( \hat{E}_g
    \hat{E}_f^* + \frac{\Omega \cdot \hat{E}_f^*}{2} \right)\hat{\sigma}_{gf}^i +
    \left( \hat{E}_f\hat{E}_g^* + \frac{\hat{E}_f \cdot \Omega}{2} \right)
    \hat{\sigma}_{fg}^i \right]
  \end{aligned}
\end{equation}

Here, we will make the assumption that the terms proportional to $\hat{\mathbf{E}}^2$
are energy shift effects that can be neglected.

\begin{equation}
  \boxed{
  \begin{aligned}
    \hat{H} = & \int d^3r \int_0^\infty \omega \hat{\mathbf{f}}^\dagger (\mathbf{r},
    \omega) \hat{\mathbf{f}}(\mathbf{r}, \omega) d\omega + \sum_{i=1}^N
    \omega_f\hat{\sigma}_{ff}^i\\
    &+\frac{1}{\Delta} \sum_{i=1}^N \left(\frac{\hat{E}_g
    \cdot \Omega^*}{2} + \frac{\Omega \cdot \hat{E}^*_g}{2} + \frac{\Omega\Omega^*}{4}\right)
    \hat{\sigma}_{gg}^i \\ &+ \frac{1}{\Delta}\sum_{i=1}^N \left[ \frac{\Omega \cdot
    \hat{E}_f^*}{2} \hat{\sigma}_{gf}^i + \frac{\hat{E}_f \cdot \Omega^*}{2}
    \hat{\sigma}_{fg}^i \right]
  \end{aligned}}
\end{equation}

\subsection{Master Equation}

In the Heisenberg picture, an operator $\hat{O}$ evolves by the equation
\begin{equation}\dot{\hat{O}} = -\frac{i}{\hbar} [\hat{O}, \hat{H}]\end{equation}
Since the operator $\hat{O}$ is atomic, we can make the observation that $[\hat{O},
\hat{\mathbf{f}}] = [\hat{O}, \hat{\mathbf{E}}] = 0$.
\begin{equation}
\begin{aligned}
  [\hat{O}, \hat{H}] = &\sum_{i=1}^N \omega_f [\hat{O}, \hat{\sigma}_{ff}^i] +
  \frac{\Omega^2}{4\Delta}[\hat{O}, \hat{\sigma}_{gg}^i] + \left( \frac{\hat{E}_g \cdot \Omega}{2\Delta} + \frac{\Omega \cdot
  \hat{E}_g^*}{2\Delta} \right)[\hat{O}, \hat{\sigma}_{gg}^i]\\ &+ \frac{\Omega \cdot
  \hat{E}_f^*}{2\Delta} [\hat{O}, \hat{\sigma}_{gf}^i]+
  \frac{\hat{E}_f \cdot \Omega}{2\Delta}
  [\hat{O}, \hat{\sigma}_{fg}^i]
\end{aligned}
\end{equation}

We can expand out the electric field in to the positive and negative field components,
and for simplicities' sake, we only look at terms proportional to the electric field
here:
\begin{align*}
  &[\mathbf{d}_{ge}e^{-i\omega_Lt}
  (\hat{\mathbf{E}}^++\hat{\mathbf{E}^-})][\hat{O}, \hat{\sigma}_{gg}^i] +
  [\mathbf{d}_{ge}^*e^{i\omega_Lt}
  (\hat{\mathbf{E}}^++\hat{\mathbf{E}^-})]
  \hat{E}_g^*[\hat{O}, \hat{\sigma}_{gg}^i]\\&+
  [\mathbf{d}_{fe}^*e^{i(\omega_L-\omega_f)t}
  (\hat{\mathbf{E}}^++\hat{\mathbf{E}^-})][\hat{O}, \hat{\sigma}_{gf}^i]+
  [\mathbf{d}_{fe}e^{-i(\omega_L-\omega_f)t}
  (\hat{\mathbf{E}}^++\hat{\mathbf{E}^-})][\hat{O}, \hat{\sigma}_{fg}^i]
\end{align*}

We assume normal ordering for the positive and negative components of the electric
field, with the positive field component acting on the commutator from the left, and the
negative field component acting on the commutator from the right.

We apply the Markov approximation for the electric field as derived in
Appendix \ref{appendix:markov}, and make the rotating wave approximation, omitting
terms that have have fast oscillations:
% RWA part 2! For the terms that are still rotating, we get rid of them.
\begin{equation}\label{Odot}
  \begin{aligned}
  \dot{\hat{O}} = &-\frac{i}{\hbar}\sum_i \omega_f [\hat{O}, \hat{\sigma}_{ff}^i] +
  \frac{\Omega^2}{4\Delta}[\hat{O}, \hat{\sigma}_{gg}^i]\\&+ \frac{\Omega}{2\Delta}
  [\hat{O}, \hat{\sigma}_{gg}^i]\left[ \frac{\mu_0\Omega}{2\Delta}\sum_j \omega_L^2
  \mathbf{G(r, r}_j, \omega_L) \mathbf{d}_{ge}\mathbf{d}_{ge}^*\hat{\sigma}_{gg}^j\;
  {\color{red}+\;
  \omega_{Lf}^2\mathbf{G(r, r}_j, \omega_{Lf})\mathbf{d}_{ge}\mathbf{d}_{fe}^*
  \hat{\sigma}_{gf}^je^{-i\omega_ft}} \right]\\&+ \frac{\Omega}{2\Delta}
  [\hat{O}, \hat{\sigma}_{gf}^i]\left[ \frac{\mu_0\Omega}{2\Delta}\sum_j {\color{red}\omega_L^2
  \mathbf{G(r, r}_j, \omega_L) \mathbf{d}_{fe}^*\mathbf{d}_{ge}\hat{\sigma}_{gg}^j
e^{-i\omega_ft}+}\; \omega_{Lf}^2\mathbf{G(r, r}_j, \omega_{Lf})\mathbf{d}_{fe}^*
  \mathbf{d}_{fe} \hat{\sigma}_{fg}^j \right] \\&+ \frac{\Omega}{2\Delta}
  \left[ \frac{\mu_0\Omega}{2\Delta}\sum_j \omega_L^2
  \mathbf{G(r, r}_j, \omega_L) \mathbf{d}_{ge}^*\mathbf{d}_{ge}\hat{\sigma}_{gg}^j
  {\color{red}\;+\;
  \omega_{Lf}^2\mathbf{G(r, r}_j, \omega_{Lf})\mathbf{d}_{ge}^*\mathbf{d}_{fe}
\hat{\sigma}_{gf}^je^{i\omega_ft}} \right][\hat{O}, \hat{\sigma}_{gg}^i]\\&+ \frac{\Omega}{2\Delta}
\left[ \frac{\mu_0\Omega}{2\Delta}\sum_j {\color{red}\omega_L^2
  \mathbf{G(r, r}_j, \omega_L) \mathbf{d}_{fe}\mathbf{d}_{ge}^*\hat{\sigma}_{gg}^j
e^{i\omega_ft} +\;}\omega_{Lf}^2\mathbf{G(r, r}_j, \omega_{Lf})\mathbf{d}_{fe}
  \mathbf{d}_{fe}^* \hat{\sigma}_{gf}^j \right][\hat{O}, \hat{\sigma}_{fg}^i]
  \end{aligned}
\end{equation}

% We define the complex coupling rate(s)
% \begin{align}
%   g^{(1)}_{ij} &= J^{(1)}_{ij} + \frac{i}{2}\Gamma^{(1)}_{ij} =
%   \frac{\mu_0\omega_L^2}{\hbar}\mathbf{d}^*_i \mathbf{G(r}_i, \mathbf{r}_j,
%   \omega_L)\mathbf{d}_j\\
%   g^{(2)}_{ij} &= J^{(2)}_{ij} + \frac{i}{2}\Gamma^{(2)}_{ij} =
%   \frac{\mu_0\omega_{Lf}^2}{\hbar}\mathbf{d}^*_i \mathbf{G(r}_i, \mathbf{r}_j,
%   \omega_{Lf})\mathbf{d}_j
% \end{align}

% How do we distinguish $\mathbf{d}_{ge}$ and $\mathbf{d}_{fe}$ if we define a general
% coupling? $g_{ij}$ could be between $ge$ and $fe^*$ or $ge$ and $ge^*$. 
Define the following couplings:
\begin{align}
  g^{(1)}_{ij} &= J^{(1)}_{ij} + \frac{i}{2}\Gamma^{(1)}_{ij} =
  \frac{\mu_0\omega_L^2}{\hbar}\mathbf{d}_{ge}^{i*} \mathbf{G(r}_i, \mathbf{r}_j,
  \omega_L)\mathbf{d}_{ge}^j\\
  g^{(2)}_{ij} &= J^{(2)}_{ij} + \frac{i}{2}\Gamma^{(2)}_{ij} =
  \frac{\mu_0\omega_{Lf}^2}{\hbar}\mathbf{d}^{i*}_{fe} \mathbf{G(r}_i, \mathbf{r}_j,
  \omega_{Lf})\mathbf{d}^j_{fe}
\end{align}

We substitute these couplings in to equation \ref{Odot}:
\begin{equation}
  \begin{aligned}
  \dot{\hat{O}} = &-\frac{i}{\hbar}\sum_i \omega_f [\hat{O}, \hat{\sigma}_{ff}^i] +
  \frac{\Omega^2}{4\Delta}[\hat{O}, \hat{\sigma}_{gg}^i]\\&+ \frac{\Omega}{2\Delta}
  [\hat{O}, \hat{\sigma}_{gg}^i]\left[ \frac{\hbar\Omega}{2\Delta}\sum_j
  g^{(1)*}_{ij}\hat{\sigma}_{gg}^j 
  \right]+ \frac{\Omega}{2\Delta}
  [\hat{O}, \hat{\sigma}_{gf}^i]\left[ \frac{\hbar\Omega}{2\Delta}\sum_j  g_{ig}^{(2)}\hat{\sigma}_{fg}^j \right] \\&+ \frac{\Omega}{2\Delta}
  \left[ \frac{\hbar\Omega}{2\Delta}\sum_j g_{ij}^{(1)}\hat{\sigma}_{gg}^j \right][\hat{O}, \hat{\sigma}_{gg}^i]+ \frac{\Omega}{2\Delta}
  \left[ \frac{\hbar\Omega}{2\Delta}\sum_j  g_{ij}^{(2)*}\hat{\sigma}_{gf}^j \right][\hat{O}, \hat{\sigma}_{fg}^i]
  \end{aligned}
\end{equation}

In the density matrix formalism, the expectation value of an operator $\hat{O}$ is
given by \[\langle \hat{O} \rangle = \text{Tr}(\hat{\rho} O)\]


We multiply both sides by $\rho(0)$ from the left and take the trace:

\begin{equation}
  \begin{aligned}
    \text{Tr}\left\{\rho(0)\dot{\hat{O}}\right\} =
    &-\frac{i}{\hbar}\text{Tr}\Bigg\{\sum_i \hbar\omega_f \rho(0)[\hat{O},
      \hat{\sigma}_{ff}^i] +
  \frac{\hbar\Omega^2}{4\Delta}\rho(0)[\hat{O}, \hat{\sigma}_{gg}^i]\\&+
  \frac{\hbar\Omega^2}{4\Delta^2} \rho(0) \sum_{i,j} \left[ g_{ij}^{(1)*}[\hat{O},
\hat{\sigma}_{gg}^i]\hat{\sigma}_{gg}^j + g_{ij}^{(1)}\hat{\sigma}_{gg}^j
[\hat{O}, \hat{\sigma}_{gg}^i]\right]\\&+
  \frac{\hbar\Omega^2}{4\Delta^2} \rho(0) \sum_{i,j} \left[ g_{ij}^{(2)}[\hat{O},
\hat{\sigma}_{gf}^i]\hat{\sigma}_{fg}^j + g_{ij}^{(2)*}\hat{\sigma}_{gf}^j
[\hat{O}, \hat{\sigma}_{fg}^i]\right] \Bigg\}\end{aligned}
\end{equation}


Using the cyclic property of the trace, we transfer the time dependence from the
operators to the density matrix to get an equation of the form
\begin{equation}
  \text{Tr}\left\{\dot{\rho(t)}\hat{O}\right\} =
  \text{Tr}\left\{\cdots\hat{O}(0)\right\}
\end{equation}

We write the master equation as

\begin{equation}
  \dot{\hat{\rho}} = \frac{i}{\hbar}[\hat{H}_0, \hat{\rho}] + \mathcal{L}[\hat{\rho}]
\end{equation}

where we define the conservative Hamiltonian
\begin{equation}
  \hat{H}_0 = \sum_{i=1}^N \hbar\omega_f\hat{\sigma}_{ff}^i + \sum_{i=1}^N
  \frac{\hbar\Omega^2}{4\Delta}\hat{\sigma}_{gg}^i + \frac{\Omega^2}{4\Delta^2}
  \sum_{i,j=1}^N \left(J_{ij}^{(1)}\hat{\sigma}_{gg}^i\hat{\sigma}_{gg}^j +
  J_{ij}^{(2)}\hat{\sigma}_{gf}^i\hat{\sigma}_{fg}^j\right)
\end{equation}

and the dissipative Linbladian
\begin{equation}
  \mathcal{L}[\hat{\rho}] = \frac{\Omega^2}{4\Delta^2} \sum_{i,j=1}^N \sum_{\alpha = \{1,2\}}
  \Gamma_{ij}^{(\alpha)} \left( \hat{L}_{\alpha}^j \hat{\rho}
    \hat{L}_{\alpha}^{i\dagger}-
  \frac{1}{2}\{\hat{L}_{\alpha}^{i\dagger}\hat{L}_{\alpha}^j,\hat{\rho}\} \right)
\end{equation}

with $\hat{L}_1 = \hat{\sigma}_{gg}$, and $\hat{L}_2 = \hat{\sigma}_{fg}$.

\newpage
\appendix
\section{Markov Approximation}\label{appendix:markov}

From \cite{PhysRevA.53.1818}, we define the electric field and the bosonic field
opearators as:
\begin{align}
\hat{\mathbf{E}}(\mathbf{r}, \omega, t) &=
i\mu_0\omega^2\sqrt{\frac{\hbar\epsilon_0}{\pi}} \int
d\mathbf{r}'\sqrt{\epsilon_{\text{I}}(\mathbf{r}',\omega)}\mathbf{G(r,r}',\omega)
\cdot\hat{\mathbf{f}}(\mathbf{r}',\omega, t)\\ \dot{\hat{\mathbf{f}}}(\mathbf{r},
\omega, t) &= \frac{i}{\hbar} \left[ \hat{H}, \hat{\mathbf{f}}(\mathbf{r}, \omega, t)
\right]
\end{align}

We recall the commutation relations
\begin{gather}
  [\hat{f}_k(\mathbf{r}, \omega), \hat{f}_{k'}^\dagger(\mathbf{r}', \omega')] =
  \delta_{kk'} \delta(\mathbf{r-r}')\delta(\omega-\omega')\\
  [\hat{f}_k(\mathbf{r}, \omega), \hat{f}_{k'}(\mathbf{r}',\omega')] = 0
\end{gather}

to solve for the bosonic field operator.

Let $H_F$ denote the field component of the Hamiltonian, and $H_{\text{int}}$ be the
rest of it.
\begin{equation}
  \dot{\hat{\mathbf{f}}}(\mathbf{r}, \omega, t) = \frac{i}{\hbar}\left( \left[
  \hat{H}_F, \hat{\mathbf{f}}\right] + \left[ \hat{H}_{\text{int}},
  \hat{\mathbf{f}}\right]\right)
\end{equation}

We solve the field commutator first
\begin{align*}
  \left[ \hat{H}_F, \hat{\mathbf{f}} \right] &= \int d^3\mathbf{r} \int_0^\infty
  d\omega\omega \left[ \hat{\mathbf{f}}^\dagger (\mathbf{r}, \omega, t)
  \hat{\mathbf{f}} (\mathbf{r}, \omega, t), \hat{\mathbf{f}}' (\mathbf{r}', \omega',
  t')\right]\\ &= \int d^3\mathbf{r} \int_0^\infty d\omega\omega \left(
  \left[ \hat{\mathbf{f}}^\dagger, \hat{\mathbf{f}}' \right]\hat{\mathbf{f}} +
  \cancel{\hat{\mathbf{f}}^\dagger \left[ \hat{\mathbf{f}}, \hat{\mathbf{f}}'
  \right]}\right)\\ &= -\omega\hat{\mathbf{f}}(\mathbf{r}, \omega, t)
\end{align*}

The interaction Hamiltonian has terms the look like either $\hat{E}_1\hat{E}_2^* \cdot
\hat{\sigma}$, $\hat{E} \cdot \hat{\sigma}$, or just simply $\hat{\sigma}$.

We start with the simplest term, $\hat{\sigma}$. The atomic operator acts separately
from the bosonic field operator, leading to
\begin{align*}
  [\hat{\sigma}, \hat{\mathbf{f}}] = 0
\end{align*}

We can next take a look at terms in the form $\hat{E}\cdot\hat{\sigma}$.
\begin{align*}
  [\hat{E}\cdot\hat{\sigma}, \hat{\mathbf{f}}] &= \cancel{\hat{E}[\hat{\sigma},
  \hat{\mathbf{f}}]} + [\hat{E}, \hat{\mathbf{f}}]\hat{\sigma}\\ &=
  \left[\hat{\mathbf{E}}(\mathbf{r}, \omega, t) \cdot \mathbf{d} e^{\mp i\omega_Lt},
  \hat{\mathbf{f}}(\mathbf{r}', \omega', t')\right]\hat{\sigma}\\ &= \left[ E^+ + E^-,
  \hat{\mathbf{f}} \right]\hat{\sigma} \cdot \mathbf{d}e^{\mp i\omega_Lt}\\ &=
  \left[i\mu_0\omega^2\sqrt{\frac{\hbar\epsilon_0}{\pi}} \int d\mathbf{r}'
  \sqrt{\epsilon_{\text{I}}(\mathbf{r}', \omega)} \mathbf{G}^*(\mathbf{r, r}', \omega)
  \cdot \hat{\mathbf{f}}^\dagger(\mathbf{r'}, \omega, t),
  \hat{\mathbf{f}}(\mathbf{r}'', \omega', t')\right] \hat{\sigma}\cdot \mathbf{d}e^{\mp
  i\omega_Lt}\\ &= -i\mu_0\omega^2 \sqrt{\frac{\hbar\epsilon_0
  \epsilon_\text{I}(\mathbf{r}', \omega) }{\pi}} \mathbf{G}^*(\mathbf{r, r'}, \omega)
  \cdot \mathbf{d}e^{\mp i\omega_Lt} \hat{\sigma}
\end{align*}

\color{red}
We finally look at terms in the form $\hat{E}_1\hat{E}_2^* \cdot \hat{\sigma}$
\begin{align*}
  [\hat{E}_1\hat{E}_2^*\cdot\hat{\sigma}, \hat{\mathbf{f}}] &= \cancel{\hat{E}_1
  \hat{E}_2^*[\hat{\sigma}, \hat{\mathbf{f}}]} + \hat{E}_1 [\hat{E}_2^*,
  \hat{\mathbf{f}}] \hat{\sigma} + [\hat{E}_1, \hat{\mathbf{f}}]
  \hat{E}_2^*\hat{\sigma}\\ &= \hat{\mathbf{E}}(\mathbf{r}, \omega, t) \left[
  \hat{\mathbf{E}}(\mathbf{r}, \omega, t), \hat{\mathbf{f}}(\mathbf{r}', \omega',
  t') \right] \mathbf{d}_1\mathbf{d}_2 \hat{\sigma} +  \left[
  \hat{\mathbf{E}}(\mathbf{r}, \omega, t), \hat{\mathbf{f}}(\mathbf{r}', \omega', t')
  \right] \hat{\mathbf{E}}(\mathbf{r}, \omega, t) \mathbf{d}_1\mathbf{d}_2
    \hat{\sigma}\\ &= -2i\mu_0\omega^2 \sqrt{\frac{\hbar\epsilon_0
  \epsilon_\text{I}(\mathbf{r}', \omega) }{\pi}} \mathbf{G}^*(\mathbf{r, r'},
  \omega) \mathbf{d}_1\mathbf{d}_2 \; \hat{\mathbf{E}}(\mathbf{r}, \omega, t) \cdot
  \hat{\sigma}
\end{align*}
\color{black}

Now, we can put together the full commutator $[H_\text{int}, \hat{\mathbf{f}}]$.
First, define $\alpha(\mathbf{r}) \equiv
i\mu_0\omega^2\sqrt{\frac{\hbar\epsilon_0\epsilon_{\text{I}}(\mathbf{r},
\omega)}{\pi}}$

\begin{equation}
\begin{aligned}
  [\hat{H}_\text{int},\hat{\mathbf{f}}] = -1* &\sum_{i=1}^N
  \color{red}
  \frac{2\alpha(\mathbf{r})}{\Delta}\mathbf{G}^*(\mathbf{r}_i, \mathbf{r}, \omega)
  \mathbf{d}_{fe}\mathbf{d}_{fe}^* \cdot \left(\hat{\mathbf{E}}(\mathbf{r}_i,\omega,t)
\cdot \hat{\sigma}_{ff}^i\right)\\ &\color{red}+
  \frac{2\alpha(\mathbf{r})}{\Delta}\mathbf{G}^*(\mathbf{r}_i, \mathbf{r}, \omega)
  \mathbf{d}_{ge}\mathbf{d}_{fe}^* \cdot \left( \hat{\mathbf{E}}(\mathbf{r}_i, \omega,
t) \cdot \hat{\sigma}_{gf}^i \right) \cdot e^{-i\omega_f t}\\&\color{red}+
  \frac{2\alpha(\mathbf{r})}{\Delta}\mathbf{G}^*(\mathbf{r}_i, \mathbf{r}, \omega)
  \mathbf{d}_{fe}\mathbf{d}_{ge}^* \cdot \left( \hat{\mathbf{E}}(\mathbf{r}_i, \omega,
t) \cdot \hat{\sigma}_{fg}^i \right)\cdot e^{i\omega_ft}\\ &\color{red}+ \frac{2\alpha(\mathbf{r})}{\Delta}
  \mathbf{G}^*(\mathbf{r}_i, \mathbf{r}, \omega) \mathbf{d}_{ge}\mathbf{d}_{ge}^*
  \cdot \left(\hat{\mathbf{E}}(\mathbf{r}_i,\omega,t) \cdot \hat{\sigma}_{gg}^i
  \right)\\ &+ \frac{\Omega\alpha(\mathbf{r})}{2\Delta} \mathbf{G}^*(\mathbf{r}_i,
  \mathbf{r}, \omega)\cdot \left(\mathbf{d}_{ge}e^{-i\omega_Lt} +
  \mathbf{d}_{ge}^*e^{i\omega_Lt}\right) \hat{\sigma}_{gg}^i \\ &+
  \frac{\Omega\alpha(\mathbf{r})}{2\Delta} \mathbf{G}^*(\mathbf{r}_i, \mathbf{r},
  \omega)\cdot \mathbf{d}_{fe}^*e^{i(\omega_L-\omega_f)t}\hat{\sigma}_{gf}^i\\ &+
  \frac{\Omega\alpha(\mathbf{r})}{2\Delta} \mathbf{G}^*(\mathbf{r}_i, \mathbf{r},
  \omega)\cdot \mathbf{d}_{fe}e^{-i(\omega_L-\omega_f)t}\hat{\sigma}_{fg}^i
\end{aligned}
\end{equation}

Putting this together with the commutator with the field hamiltonian, the bosonic
field operator is now defined:
\begin{equation}\boxed{
\begin{aligned}
\dot{\hat{\mathbf{f}}}(\mathbf{r}, \omega, t) = &-i\omega\hat{\mathbf{f}}(\mathbf{r}, \omega,
  t)\\ &+ \sum_{i=1}^N \frac{\Omega}{2\Delta}
  \frac{\omega^2}{c^2} \sqrt{\frac{\epsilon_\text{I}(\mathbf{r,
  \omega})}{\hbar\pi\epsilon_0}} \mathbf{G}^*(\mathbf{r}_i, \mathbf{r}, \omega) \cdot
  \left( \left(\mathbf{d}_{ge}e^{-i\omega_Lt} +
  \mathbf{d}_{ge}^*e^{i\omega_Lt}\right) \hat{\sigma}_{gg}^i
  \right.\\&\hspace{10em}+\left. \mathbf{d}_{fe}^*e^{i(\omega_L-\omega_f)t}\hat{\sigma}_{gf}^i+
  \mathbf{d}_{fe}e^{-i(\omega_L-\omega_f)t}\hat{\sigma}_{fg}^i\right)\\&\color{red}+ \sum_{i=1}^N \frac{2
  \omega^2}{\Delta c^2} \sqrt{\frac{\epsilon_\text{I}(\mathbf{r},
  \omega)}{\hbar\pi\epsilon_0}} \mathbf{G}^*(\mathbf{r}_i, \mathbf{r},
  \omega)\hat{\mathbf{E}}(\mathbf{r}_i, \omega, t) \cdot \left( \mathbf{d}_{fe}
  \mathbf{d}_{fe}^*\hat\sigma_{ff}^i + \mathbf{d}_{ge}
    \mathbf{d}_{ge}^*\hat\sigma_{gg}^i \right.\\&\color{red}\hspace{13em}+ \left.\mathbf{d}_{ge}
    \mathbf{d}_{fe}^*e^{-i\omega_ft}\hat\sigma_{gf}^i + \mathbf{d}_{fe}
  \mathbf{d}_{ge}^*e^{i\omega_ft}\hat\sigma_{fg}^i \right)
\end{aligned}}
\end{equation}

This is somewhat similar to the case derived for the 2 level case in
\cite{PhysRevA.66.063810}, but now we have an extra term at the end proportional to
$\hat{\mathbf{E}}\cdot\hat{\sigma}$.\\\vspace{1em}


We can take the time derivative of the electric field operator
\begin{align*}
  \dot{\hat{\mathbf{E}}} &= -i\omega\hat{\mathbf{E}} + i\mu_0\omega^2\sqrt{\frac{\hbar\epsilon_0}{\pi}} \int
  d\mathbf{r}'\sqrt{\epsilon_{\text{I}}(\mathbf{r}',\omega)}\mathbf{G(r,r}',\omega)
  \cdot\dot{\hat{\mathbf{f}}}(\mathbf{r}',\omega, t)\\&= -i\omega\hat{\mathbf{E}}+
  i\mu_0
  \frac{\omega^4}{\pi c^2}\frac{\Omega}{2\Delta}\int d\mathbf{r}'
  \epsilon_I(\mathbf{r}',\omega)
  \mathbf{G}(\mathbf{r, r'}, \omega) \times \sum_{j=1}^N \mathbf{G}^*(\mathbf{r}_j,
  \mathbf{r}, \omega)\cdot \hat{\mathbf{d}}^{(1)}_j(t) \\&\color{red}\hspace{1em}+ i\mu_0
  \frac{\omega^4}{\pi c^2}\frac{2}{\Delta} \int d\mathbf{r}'\epsilon_\text{I}(\mathbf{r}',
  \omega)\mathbf{G}(\mathbf{r, r}', \omega) \times \sum_{j=1}^N
  \mathbf{G}^*(\mathbf{r}_j, \mathbf{r}, \omega)\cdot \hat{\mathbf{E}}(\mathbf{r}_j,
  \omega, t) \cdot \hat{\mathbf{d}}^{(2)}_{j}(t)
\end{align*}

where the dipole opearators are defined
\begin{align}
  \hat{\mathbf{d}}^{(1)}_j(t) &= \left(\mathbf{d}_{ge}e^{-i\omega_Lt} +
  \mathbf{d}_{ge}^*e^{i\omega_Lt}\right) \hat{\sigma}_{gg}^j
  + \mathbf{d}_{fe}^*e^{i(\omega_L-\omega_f)t}\hat{\sigma}_{gf}^j+
  \mathbf{d}_{fe}e^{-i(\omega_L-\omega_f)t}\hat{\sigma}_{fg}^j\\\color{red}
  \hat{\mathbf{d}}^{(2)}_j(t) &\color{red}= \mathbf{d}_{fe}
  \mathbf{d}_{fe}^*\hat\sigma_{ff}^j + \mathbf{d}_{ge}
  \mathbf{d}_{ge}^*\hat\sigma_{gg}^j +\mathbf{d}_{ge}
  \mathbf{d}_{fe}^*e^{-i\omega_ft}\hat\sigma_{gf}^j + \mathbf{d}_{fe}
  \mathbf{d}_{ge}^*e^{i\omega_ft}\hat\sigma_{fg}^j
\end{align}

We apply the following Green's function identity
\begin{equation}
  \frac{\omega^2}{c^2}\int d^3\mathbf{s}\; \epsilon_{\text{I}}(\mathbf{s}, \omega)
  G_{ik}(\mathbf{r,s}, \omega)G_{jk}^* (\mathbf{r',s}, \omega) = \text{Im}
  \left\{G_{ij}(\mathbf{r,r'}, \omega)\right\}
\end{equation}

This results in our 
\begin{equation}\label{Edot}
\begin{aligned}
  \dot{\hat{\mathbf{E}}} = &-i\omega\hat{\mathbf{E}}+
  \frac{i\mu_0\omega^2}{\pi} \frac{\Omega}{2\Delta} \sum_{j=1}^N \text{Im}\left\{
  \mathbf{G}(\mathbf{r, r}_j, \omega) \right\} \cdot \hat{\mathbf{d}}^{(1)}_j(t)\\ &\color{red}+
  \frac{i\mu_0\omega^2}{\pi}\frac{2}{\Delta} \sum_{j=1}^N \text{Im}\{ \mathbf{G(r,
  r}_j, \omega) \} \cdot \hat{\mathbf{E}}(\mathbf{r}_j, \omega, t)\cdot
  \hat{\mathbf{d}}_{j}^{(2)}(t)
\end{aligned}
\end{equation}

This can be simplified by collecting a couple terms:
\begin{equation}\label{Edott}\boxed{
\begin{aligned}
  \dot{\hat{\mathbf{E}}} = -i\omega\hat{\mathbf{E}}(\mathbf{r},\omega, t)+\frac{i\mu_0\omega^2}{\Delta\pi}\sum_{j=1}^N \text{Im}\{ \mathbf{G(r, r}_j,
  \omega) \} \cdot \left( \frac{\Omega}{2}\hat{\mathbf{d}}_j^{(1)}(t)\color{red} + 2
  \hat{\mathbf{E}}(\mathbf{r}_j, \omega, t)\cdot
\hat{\mathbf{d}}^{(2)}_j(t)\color{black} \right)
\end{aligned}}
\end{equation}

Up until this point, I have included the terms that arise from the
$\hat{\mathbf{E}}^2$ component of the eliminated Hamiltonian. As previously mentioned,
we predict the $\hat{\mathbf{E}}^2$ terms will only show up as an energy shift, which
will already be accounted for in the final equation, meaning that this will lead to a
possible double counting of said energy shift. Therefore, from here onwards, these
terms will be neglected.

\vspace{1em} We define a slowly varying variable by $\hat{\mathbf{E}}(\mathbf{r},
\omega) =
e^{-i\omega t}\tilde{\mathbf{E}}(\mathbf{r}, \omega)$. Substituting into
\ref{Edott} gives
\begin{equation}
  \dot{\tilde{\mathbf{E}}}(\mathbf{r}, \omega) =
  \frac{i\mu_0\omega^2}{\Delta\pi}\sum_{j=1}^N \text{Im}\{ \mathbf{G(r, r}_j,
    \omega) \} \cdot \frac{\Omega}{2}e^{i\omega t}\hat{\mathbf{d}}_j^{(1)}(t)
\end{equation}

Integrating this gives
\begin{equation}
  \tilde{\mathbf{E}}(\mathbf{r}, \omega, t) = \tilde{\mathbf{E}}(\mathbf{r}, \omega,
  0) + \frac{i\mu_0\omega^2}{\Delta\pi}\sum_{j=1}^N \text{Im}\left\{\mathbf{G(r, r}_j,
  \omega)\right\} \cdot \int_0^tdt'
  \frac{\Omega}{2}\hat{\mathbf{d}}_j^{(1)}(t')e^{i\omega t'}
\end{equation}
We multiply both sides by $e^{-i\omega t}$ and convert back to the normal
variables.
\begin{equation}
  \mathbf{E}(\mathbf{r}, \omega, t) = \mathbf{E}(\mathbf{r}, \omega,
  0)e^{-i\omega t} + \frac{i\mu_0\omega^2}{\Delta\pi}\sum_{j=1}^N \text{Im}\left\{\mathbf{G(r, r}_j,
  \omega)\right\} \cdot \int_0^tdt'
  \frac{\Omega}{2}\hat{\mathbf{d}}_j^{(1)}(t')e^{-i\omega (t-t')}
\end{equation}

We substitute this expression into the field operator $\hat{\mathbf{E}}^+(\mathbf{r},
t) = \int_0^\infty d\omega\hat{\mathbf{E}}^+(\mathbf{r}, \omega)$ to get
\begin{equation}
  \mathbf{E}^+(\mathbf{r}, t) = \mathbf{E}_0 +
  \sum_{j=1}^N \int_0^\infty d\omega\,\frac{i\mu_0\omega^2}{\Delta\pi}
  \text{Im}\left\{\mathbf{G(r, r}_j,
  \omega)\right\} \cdot \int_0^tdt'\,
  \frac{\Omega}{2}\hat{\mathbf{d}}_j^{(1)}(t')e^{-i\omega (t-t')}
\end{equation}

where we define $\mathbf{E}_0 = \int_0^\infty \hat{\mathbf{E}}(\mathbf{r}, \omega,
0)e^{-i\omega t}d\omega$.

We isolate the time integral
\begin{align*}
  &\frac{\Omega}{2}\int_0^tdt' \left[\left(\mathbf{d}_{ge}e^{-i\omega_Lt'} +
  \mathbf{d}_{ge}^*e^{i\omega_Lt'}\right) \hat{\sigma}_{gg}^j(t')
  + \mathbf{d}_{fe}^*e^{i(\omega_L-\omega_f)t'}\hat{\sigma}_{gf}^j(t')+
\mathbf{d}_{fe}e^{-i(\omega_L-\omega_f)t'}\hat{\sigma}_{fg}^j(t')\right]e^{-i\omega(t-t')}\\
  &= \frac{\Omega}{2}\int_0^tdt'
  \mathbf{d}_{ge}e^{-i(\omega-\omega_L)(t-t')}e^{-i\omega_Lt}\hat{\sigma}_{gg}^j(t') +
  \mathbf{d}_{ge}^* e^{-i(\omega+\omega_L)(t-t')}e^{i\omega_Lt}\hat{\sigma}_{gg}^j(t')
  \\&\hspace{4em}+
  \mathbf{d}_{fe}e^{-i(\omega -
  \omega_L+\omega_f)(t-t')}e^{-i(\omega_L-\omega_f)}\hat{\sigma}_{fg}^j(t') +
  \mathbf{d}_{fe}^*e^{-i(\omega +
  \omega_L-\omega_f)(t-t')}e^{i(\omega_L-\omega_f)}\hat{\sigma}_{gf}^j(t')
\end{align*}

We take the atomic operators to be slowly varying via the Markov approximation and
move them out of the time integral.
\begin{equation}
  \begin{aligned}
  &\hat{\sigma}_{gg}^j(t)e^{-i\omega_Lt}\int_0^tdt'
  \mathbf{d}_{ge}e^{-i(\omega-\omega_L)(t-t')} + \hat{\sigma}_{gg}^j(t)e^{i\omega_Lt}\int_0^\infty
  \mathbf{d}_{ge}^* e^{-i(\omega+\omega_L)(t-t')}
  \\&+
  \hat{\sigma}_{fg}^j(t)e^{-i(\omega_L-\omega_f)t}\int_0^\infty \mathbf{d}_{fe}e^{-i(\omega -
  \omega_L+\omega_f)(t-t')}+ \hat{\sigma}_{gf}^j(t)e^{i(\omega_L-\omega_f)t}\int_0^\infty
  \mathbf{d}_{fe}^*e^{-i(\omega +
  \omega_L-\omega_f)(t-t')}
  \end{aligned}
\end{equation}

We use the following relation
\begin{equation}
\int_{0}^tdt'\,e^{-i(\omega-\omega')(t-t')}=\zeta(\omega'-\omega)
\end{equation}

resulting in the following:
\begin{equation}
  \begin{aligned}
    \hat{\mathbf{E}}^+(\mathbf{r}, t) &= \mathbf{E}_0 + \sum_{j=1}^N\int_0^\infty
    d\omega\,\frac{i\mu_0\omega^2}{\pi}\frac{\Omega}{2\Delta}\text{Im}
    \left\{\mathbf{G(r, r}_j, \omega)\right\} \cdot \\&\bigg(
      \mathbf{d}_{ge}\hat{\sigma}_{gg}^j(t)e^{-i\omega_Lt}\zeta(\omega_L-\omega) +
      \mathbf{d}_{ge}^*\hat{\sigma}_{gg}^j(t)e^{i\omega_Lt}\zeta(\omega_L+\omega)\\& +
      \mathbf{d}_{fe}\hat{\sigma}_{fg}^j(t)e^{-i(\omega_L-\omega_f)t}\zeta(\omega_L-\omega_f-\omega) +
    \mathbf{d}_{fe}^*\hat{\sigma}_{gf}^j(t)e^{i(\omega_L-\omega_f)t}\zeta(\omega_L-\omega_f+\omega)\bigg)
  \end{aligned}
\end{equation}

We can get rid of the delta function in $\zeta(\omega_L+\omega)$ term since it doesn't
live in the integral. The principal
value integral is related to Casimir forces \cite{Chang_2014} and is small when far
away from the dielectric \cite{Hung_2013}.

If $\omega_L>\omega_f$, the $\sigma_{fg}$ term is allowed and the $\sigma_{gf}$ term
is not allowed and vice versa if $\omega_L<\omega_f$. Here, we assume
$\omega_L\neq\omega_f$. We choose $\omega_L>\omega_f$.
\begin{equation}\label{a20}
  \begin{aligned}
    \hat{\mathbf{E}}^+(\mathbf{r}, t) &= \mathbf{E}_0 + \sum_{j=1}^N\int_0^\infty
    d\omega\,\frac{i\mu_0\omega^2}{\pi}\frac{\Omega}{2\Delta}\text{Im}
    \left\{\mathbf{G(r, r}_j, \omega)\right\} \cdot \left[\mathbf{d}_{ge}\hat{\sigma}_{gg}^j(t)\zeta(\omega_L-\omega) +
    \mathbf{d}_{fe}\hat{\sigma}_{fg}^j(t)\zeta(\omega_L-\omega_f-\omega)\right]
  \end{aligned}
\end{equation}

where we have redefined $\hat{\sigma}_{gg}^j(t) =
\hat{\sigma}_{gg}^j(t)e^{-i\omega_Lt}$ and $\hat{\sigma}_{fg}^j(t) =
\hat{\sigma}_{fg}^j(t)e^{-i(\omega_L-\omega_f)t}$.

We now expand the function zeta
\begin{equation}
  \zeta(x) = \pi\delta(x) + \mathcal{P}\frac{i}{x}
\end{equation}
and insert it into \ref{a20}:
\begin{align*}
  \hat{\mathbf{E}}^+(\mathbf{r}, t) &= \mathbf{E}_0 + \sum_{j=1}^N \int_0^\infty
  d\omega\, \frac{i\mu_0\omega^2}{\pi}\frac{\Omega}{2\Delta} \text{Im}\left\{
  \mathbf{G(r, r}_j,\omega) \right\}\mathbf{d}_{ge}\hat{\sigma}_{gg}^j(t) \left(
\pi\delta(\omega_L - \omega) + \mathcal{P}\frac{i}{\omega_L - \omega} \right)\\
                                    &\hspace{2.5em}+
\sum_{j=1}^N \int_0^\infty d\omega \frac{i\mu_0\omega^2}{\pi} \frac{\Omega}{2\Delta}
\text{Im}\left\{ \mathbf{G(r, r}_j, \omega) \right\}
\mathbf{d}_{fe}\hat{\sigma}_{fg}^j(t) \left( \pi\delta(\omega_L-\omega_f-\omega) +
\mathcal{P}\frac{i}{\omega_L-\omega_f-\omega} \right)\\ &= \mathbf{E}_0 + \sum_{j=1}^N
\mathbf{d}_{ge}\hat{\sigma}_{gg}^j \left[ \frac{i\mu_0\omega_L^2\Omega}{2\Delta}
\text{Im} \left\{ \mathbf{G(r, r}_j, \omega_L) \right\} -
\frac{\mu_0\Omega}{2\pi\Delta} \mathcal{P} \int_0^\infty d\omega \frac{\omega^2
\text{Im}\left\{ \mathbf{G(r, r}_j, \omega) \right\}}{\omega_L-\omega} \right]\\
                                                        &\hspace{2.5em}+ \sum_{j=1}^N
\mathbf{d}_{fe}\hat{\sigma}_{fg}^j \left[ \frac{i\mu_0(\omega_L-\omega_f)^2\Omega}{2\Delta}
\text{Im} \left\{ \mathbf{G(r, r}_j, \omega_L - \omega_f) \right\} -
\frac{\mu_0\Omega}{2\pi\Delta} \mathcal{P} \int_0^\infty d\omega \frac{\omega^2
\text{Im}\left\{ \mathbf{G(r, r}_j, \omega) \right\}}{\omega_L - \omega_f -\omega} \right]
\end{align*}

The greens function is analytic in the upper half of the complex plane, and the
Kramers-Kronig relation is
\begin{equation}
  \text{Re}\left\{ \mathbf{G(r, r}_j, \omega_L) \right\} = \frac{1}{\pi}\mathcal{P}
  \int_{-\infty}^\infty \frac{\text{Im}\left\{\mathbf{G(r, r}_j,
  \omega)\right\}}{\omega - \omega_L}
 = \frac{2}{\pi}\mathcal{P}
  \int_0^\infty \frac{\omega \text{Im}\left\{\mathbf{G(r, r}_j,
  \omega)\right\}}{\omega - \omega_L}
\end{equation}
by use of the property that the real part of the greens function is odd in omega.

The principal value integral in our case has an extra $\omega$, and the correction
term for that is derived in \cite{Dzsotjan_2011}:
\begin{equation}
  \mathcal{P}\int_0^\infty d\omega \frac{\omega^2 \text{Im}\left\{\mathbf{G(r, r}_j,
  \omega)\right\}}{\omega-\omega_A} = \pi\omega_A^2 \text{Re}\left\{\mathbf{G(r, r}_j,
  \omega_A)\right\} + \int_0^\infty d\kappa \kappa^2 \text{Re}\left\{\mathbf{G(r, r}_j,
  \kappa)\right\} \frac{\omega_A}{\kappa^2+\omega_A^2}
\end{equation}

This integral term is negligible for separation distances much larger than the
characteristic length \cite{Dzsotjan_2011}.

Expanding out the integral, and combining the real and imaginary parts of the Greens
function, we get
\begin{equation}\boxed{
  \hat{\mathbf{E}}^+(\mathbf{r}, t) = \mathbf{E}_0 + \frac{\mu_0\Omega}{2\Delta}
  \sum_{j=1}^N \omega_L^2\mathbf{G(r, r}_j, \omega_L) \cdot
  \mathbf{d}_{ge} \hat{\sigma}_{gg}^je^{-i\omega_Lt} + \omega_{Lf}^2
  \mathbf{G(r, r}_j, \omega_{Lf}) \cdot \mathbf{d}_{fe}
\hat{\sigma}_{fg}^je^{-i(\omega_L-\omega_f)t}
}\end{equation}
where $\omega_{Lf} = \omega_L-\omega_f$.

\color{black}
\newpage
\bibliography{refs}{}
\bibliographystyle{plain}
\end{document}
